% p.27 — الباب الرابع العشر في دعا يدعو (تابع): المنفيات والمؤكدات
\mspage{27}
% Layout: six tall columns (right-to-left), each a stack of boxed blocks of unequal height.
% Every sub-heading box holds the heading AND, beneath it, the underlying (pre-iʿlāl) form,
% e.g. «مجهول بلم / لَمْ يُدْعَوْ»; the regular forms follow outside the box. Row alignment is
% approximate: all six columns run to 21 lines, so rows are aligned 1:1; the \cline marks
% reproduce the printed box boundaries inside each column.
% Column 1 continues the مضارع بلم (معروف) block begun on p.26. In the ناقص واوي the underlying
% and the actual forms of the معروف (e.g. لَا يَدْعُوَنَّ) coincide, so the lithograph prints them twice.
% The running header line of the frame is blank (smudged) on this page.
\begin{sarf}{|c|c|c|c|c|c|}
لَمْ تَدْعُوَا & لَمْ أُدْعَ & لَنْ يُدْعَيَا & لَتَدْعُوَنَّ & لَتُدْعَيْنَانِّ & \thead{أمر حاضر بالنون}\\
لَمْ تَدْعُوْا & لَمْ نُدْعَ & لَنْ يُدْعَوْا & لَتَدْعُوَانِّ & لَأُدْعَيَنَّ & اُدْعُوَنَّ\\\cline{2-2}\cline{6-6}
لَمْ تَدْعِيْ & \thead{مضارع بلن} & لَنْ تُدْعَى & لَتَدْعُنَّ & لَنُدْعَيَنَّ & اُدْعُوَنَّ\\\cline{5-5}
لَمْ تَدْعُوَا & لَنْ يَدْعُوَ & لَنْ تُدْعَيَا & لَتَدْعِنَّ & \thead{منفي بالنون} & اُدْعُوَانِّ\\\cline{2-2}
لَمْ تَدْعُوْنَ & لَنْ يَدْعُوَ & لَنْ يُدْعَيْنَ & لَتَدْعُوَانِّ & لَا يَدْعُوَنَّ & اُدْعُنَّ\\\cline{5-5}
لَمْ أَدْعُ & لَنْ يَدْعُوَا & لَنْ تُدْعَى & لَتَدْعُوْنَانِّ & لَا يَدْعُوَنَّ & اُدْعِنَّ\\\cline{5-5}
لَمْ نَدْعُ & لَنْ يَدْعُوْا & لَنْ تُدْعَيَا & \com{لَا دْعُوَنَّ}{كذا في الأصل، بألف بلا همزة؛ ولعل الصواب: لَأَدْعُوَنَّ} & \thead{مجهول بالنون} & اُدْعُوَانِّ\\\cline{1-1}
\thead{مجهول بلم} & لَنْ تَدْعُوَ & لَنْ تُدْعَوْا & لَنَدْعُوَنَّ & لَا يُدْعَوَنَّ & اُدْعُوْنَانِّ\\\cline{4-4}\cline{5-5}\cline{6-6}
لَمْ يُدْعَوْ & لَنْ تَدْعُوَا & لَنْ تُدْعَيْ & \thead{مجهول بالنون} & لَا يُدْعَيَنَّ & \thead{نهي حاضر بالنون}\\\cline{1-1}\cline{5-5}
لَمْ يُدْعَ & لَنْ يَدْعُوْنَ & لَنْ تُدْعَيَا & لَيُدْعَوَنَّ & \thead{أمر غائب بالنون} & لَا تَدْعُوَنَّ\\\cline{4-4}\cline{6-6}
لَمْ يُدْعَيَا & لَنْ تَدْعُوَ & لَنْ تُدْعَيْنَ & لَيُدْعَيَنَّ & لِيَدْعُوَنَّ & لَا تَدْعُوَنَّ\\\cline{5-5}
لَمْ يُدْعَوْا & لَنْ تَدْعُوَا & لَنْ أُدْعَى & لَيُدْعَيَانِّ & لِيَدْعُوَنَّ & لَا تَدْعُوَانِّ\\\cline{5-5}
لَمْ تُدْعَ & لَنْ تَدْعُوْا & لَنْ نُدْعَى & لَيُدْعَوُنَّ & \thead{مجهول بالنون} & لَا تَدْعُنَّ\\\cline{3-3}
لَمْ تُدْعَيَا & لَنْ تَدْعِيْ & \thead{مضارع بالنون} & لَتُدْعَيَنَّ & لِيُدْعَوَنَّ & لَا تَدْعِنَّ\\\cline{5-5}
لَمْ يُدْعَيْنَ & لَنْ تَدْعُوَا & لَيَدْعُوَنَّ & لَتُدْعَيَانِّ & لِيُدْعَيَنَّ & لَا تَدْعُوَانِّ\\\cline{3-3}\cline{5-5}
لَمْ تُدْعَ & لَنْ تَدْعُوْنَ & لَيَدْعُوَنَّ & لَيُدْعَيْنَانِّ & \thead{نهي غائب بالنون} & لَا تَدْعُوْنَانِّ\\\cline{6-6}
لَمْ تُدْعَيَا & لَنْ أَدْعُوَ & لَيَدْعُوَانِّ & لَتُدْعَيَنَّ & لَا يَدْعُوَنَّ & \thead{تمرين}\\\cline{5-5}\cline{6-6}
لَمْ تُدْعَوْا & لَنْ نَدْعُوَ & لَيَدْعُنَّ & لَتُدْعَيَانِّ & لَا يَدْعُوَنَّ & غَزَا\\\cline{2-2}\cline{5-5}
لَمْ تُدْعَيْ & \thead{مجهول بلن} & لَتَدْعُوَنَّ & لَتُدْعَوُنَّ & \thead{مجهول بالنون} & صَفَا\\
لَمْ تُدْعَيَا & لَنْ يُدْعَوَ & لَتَدْعُوَانِّ & لَتُدْعَيِنَّ & لَا يُدْعَوَنَّ & طَفَا\\\cline{2-2}\cline{5-5}
لَمْ تُدْعَيْنَ & لَنْ يُدْعَى & لَيَدْعُوْنَانِّ & لَتُدْعَيَانِّ & لَا يُدْعَيَنَّ & عَزَا\\
\end{sarf}
